\begin{abstract}
The choice of which problems to train on during code SFT matters far more than practitioners
assume. We show that training DeepSeek-Coder-1.3B on HumanEval algorithm problems versus MBPP
utility scripts versus LeetCode challenges --- under identical token budgets and controlled
conditions --- produces HumanEval+ pass@1 differences of up to 31.9 percentage points (ANOVA
F=11.37, p=0.020). More surprisingly, the model trained on HumanEval problems achieves the
best MBPP+ performance too ($\sim$52\% vs $\sim$50.5\% for MBPP-only SFT), violating the intuitive
``train on X to test on X'' principle. We show that this is not coincidental: CodeBERT
embedding cosine similarity between the training source and the test benchmark perfectly
predicts the pass@1 rank across all four conditions (Spearman $\rho$=1.0, p=0.042,
dual-encoder concordant) --- the first permutation-tested evidence that distributional alignment
in code-embedding space governs SFT specialization. These findings suggest that embedding
similarity to the target benchmark is a principled, pre-training predictor of SFT performance,
and that algorithmic function-completion training develops more transferable Python programming
competence than utility-script training.
\end{abstract}
